% ---------------------------------------------------------------------------
% Cleaved HA (HA1 | HA2) in ESM3: what was built, what was measured, what is
% still open
%
% Session report for the bridge_test investigation of 2026-10-08/09.
% Companion to softmax/code/bridge_test/DEVLOG_bridge_test.txt and
% DEVLOG_calibration.txt: those are the chronological record, this is the
% standing summary with the limitations and the to-do list.
%
% Build:  cd softmax/docs && pdflatex bridge_test_report.tex && pdflatex bridge_test_report.tex
% (twice, for the table of contents and cross-references)
% ---------------------------------------------------------------------------
\documentclass[11pt,a4paper]{article}

\usepackage[T1]{fontenc}
\usepackage[utf8]{inputenc}
\usepackage[margin=2.5cm]{geometry}
\usepackage{amsmath,amssymb}
\usepackage{booktabs}
\usepackage{array}
\usepackage[hidelinks]{hyperref}
% File paths are set with hyperref's \path, which breaks at slashes and needs
% no underscore escaping; \sloppy absorbs the rest.
\sloppy

% \mbox keeps \AA in text mode, so \ang works both inside and outside $...$.
\newcommand{\ang}{\,\mbox{\AA}}
\newcommand{\ptm}{\mathrm{pTM}}
\newcommand{\plddt}{pLDDT}

\title{Cleaved hemagglutinin in ESM3:\\
       \large Can the model treat HA1 and HA2 as one molecule\\
       joined by a disulfide bridge?}
\author{\texttt{my\_extendedSampling} / \texttt{softmax} / \texttt{bridge\_test}}
\date{2026--10--09}

\begin{document}
\maketitle

\begin{abstract}
\noindent
\texttt{zero\_polymer} is an influenza hemagglutinin precursor (HA0). In the
mature protein it is cut once, after Arg329 in mature numbering, into HA1 and
HA2, which stay together through an interchain disulfide between Cys14 of HA1
and Cys137 of HA2. The question was whether ESM3 can be given that two-chain
molecule and handle it. The mechanical answer is yes: the stock tokenizer
accepts the chain-break character, but no benefit of declaring the break shows
up consistently across decoding step counts. The structural answer is no, with
caveats: no
assembled construct folds well enough to pass the geometry gate, and the
interchain Cys14--Cys137 distance is $96\ang$ (cleaved) and $130\ang$ (no
break), against a bonding criterion of $4.5\ang$. Experimental structures of the two disulfide controls then showed that the
$4.5\ang$ criterion is exact on real geometry: the native bonds, and nothing
else, lie within it. The failures are in the predictions. The predicted
crambin fold is wrong (TM-score $0.32$) at pLDDT $83$ and passes the geometry
gate, and the model places a non-native pair at $3.3\ang$ that is $17.9\ang$ in
the experiment. The round trip reproduces both folds to about $0.3\ang$, so the
representation can hold the bonds; a coordinate prompt pulls a pair partway
toward native geometry without holding it. This report records the results,
the corrections made along the way, the limitations, and a prioritised to-do
list.
\end{abstract}

\tableofcontents

% ===========================================================================
\section{Summary}
% ===========================================================================

\begin{itemize}
\item \textbf{Mechanics work.} \texttt{model.encode(ESMProtein(sequence=...))}
  accepts \texttt{|}; 551 positions go in and 551 come out;
  \texttt{to\_protein\_complex()} recovers two chains. No change to the
  model code was needed.
\item \textbf{No consistent benefit of the chain-break.} At 256 steps HA1's
  mean \plddt{} is 51.4 fused to HA2, 63.2 with the break and 65.2 alone, which
  looked like a benefit. Across 64, 128 and 256 steps the break-minus-fused
  difference is $-5.2$, $+2.4$, $+11.8$: the sign is not stable, and the
  earlier claims that the break helps HA1 and its Cys14 anchor are withdrawn.
\item \textbf{Step count is a large, non-monotone effect.} HA1 alone scores
  48.4, 67.0, 65.2 and 59.8 at 64, 128, 256 and 329 steps. Each step count is a
  deterministic but effectively independent draw, and draws differ by more
  than most effects claimed so far. Only sign-stable differences are findings.
\item \textbf{HA2 alone folds well; every assembly degrades it.} HA2 alone
  passes the gate at all three step counts (pLDDT 88--94, no clashes) and is
  17--29 points more confident than HA2 in any assembly. No assembled
  construct, and no HA1-alone run, passes the gate at any step count.
\item \textbf{The interchain bridge is not formed} in either assembled
  construct, by a margin (at least $10\times$ the criterion at every step count) that no
  threshold subtlety can change.
\item \textbf{The $4.5\ang$ criterion is exact on experimental structures.} In
  BPTI and crambin, exactly the native bonds lie within $4.5\ang$ and nothing
  else does. The earlier verdict that it lacks specificity was wrong: the
  non-native pair is a prediction error.
\item \textbf{The predicted crambin fold is wrong} (TM-score $0.32$), at pLDDT
  83 and with no clashes. The geometry gate is necessary, not sufficient.
\item \textbf{The representation can hold the bonds.} The round trip reproduces
  both folds to about $0.3\ang$. A coordinate prompt pulls a bridge partway
  toward native geometry, locally, at a cost in confidence, and does not hold it.
\item \textbf{Not yet tested:} the controls at other step counts, the trimer, an experimental HA structure, partner identity for HA.
\end{itemize}

% ===========================================================================
\section{What was built}
% ===========================================================================

All of it is on branch \path{claude/esm3-disulfide-polymer-chains-rwvpoq}.

\begin{center}
\small
\begin{tabular}{@{}p{5.6cm}p{9.8cm}@{}}
\toprule
Commit & Content \\
\midrule
\texttt{46bccd8} & \path{tests/references.py}: chains, bridge, numbering helpers \\
\texttt{1f35fbf}, \texttt{fcb86aa}, \texttt{caeb436} &
  \path{tests/test_cleaved_complex.py}, three revisions \\
\texttt{4056d91}, \texttt{346b3de} &
  BPTI and crambin controls; known-disulfide measurement \\
\texttt{086e2cd}, \texttt{1b0d3ce}, \texttt{6b34313}, \texttt{8efcd2c} &
  \path{bridge_test/DEVLOG_bridge_test.txt} \\
\texttt{87dc571} & \path{bridge_test/DEVLOG_calibration.txt} \\
\bottomrule
\end{tabular}
\end{center}

\paragraph{The construct.} Chain one is mature HA1 (329 aa, ends
\texttt{\ldots QTR}); chain two is HA2 (221 aa, starts \texttt{GIFGAI}); the
16-residue signal peptide is dropped, so 550 residues in total. Both chains
are slices of \texttt{ZERO\_POLYMER}, never retyped, and \texttt{\_check()}
asserts the junction, the lengths, and that both anchors are cysteines.

\paragraph{A numbering trap worth keeping.} The HA landmarks are in
\emph{mature} numbering, which runs 16 behind the index into
\texttt{ZERO\_POLYMER}. Read naively, \texttt{ZERO\_POLYMER[328]} (``Arg329'')
is \texttt{T} and \texttt{ZERO\_POLYMER[13]} (``Cys14'') is \texttt{V}.
Neither raises. \texttt{mature\_to\_index()} and \texttt{chain\_index()} exist
so the conversion is never done by hand.

% ===========================================================================
\section{Method}
% ===========================================================================

\paragraph{Route into the model.} The sampler's usual path
(\texttt{ExtendedProtein} with \texttt{predict\_attention} on soft
\texttt{sequence\_probs}) cannot carry a chain break: \texttt{expand()} maps
through a 25-letter vocabulary with no \texttt{|}, and the encoder zero-pads
over token 31. The stock \texttt{encode}/\texttt{predict\_protein} route runs
the real tokenizer and has none of these problems. Two shortcuts that look
right are wrong: \texttt{chain\_id=} is consumed only by geometric attention,
which the soft path masks and zeroes, so it is a silent no-op; and
\texttt{sequence\_id=} is an attention separator, so distinct ids would make
the chains ignore each other.

\paragraph{What the decoder returns.} Backbone plus an inferred C$\beta$, and
no side-chain sulfur at any residue. So SG--SG is unavailable and
C$\beta$--C$\beta$ distance is the only disulfide proxy. C$\beta$ is not
predicted: \texttt{infer\_cbeta()} derives it from N/CA/C and writes NaN at
every glycine. The chain-break position is excluded from all C$\beta$
statistics.

\paragraph{The instrument.}
\begin{itemize}
\item \emph{Geometry gate:} pTM $\ge 0.40$, and no C$\beta$--C$\beta$ pair
  closer than $3.0\ang$ at sequence separation $\ge 2$.
\item \emph{Bond criterion:} C$\beta$--C$\beta$ $\le 4.5\ang$.
\item \emph{Controls:} protein~G (fold control, no cysteines); BPTI and
  crambin (three known disulfides each).
\end{itemize}

% ===========================================================================
\section{Results}
% ===========================================================================

\subsection{The controls}

\begin{center}
\begin{tabular}{@{}lrrrrrr@{}}
\toprule
Reference & steps & pTM & pLDDT & \% $\ge 70$ & min C$\beta$ ($\ang$) & clashes \\
\midrule
protein G (56 aa) & 56 & 0.7596 & 93.24 & 100.0 & 3.68 & 0 \\
BPTI (58 aa)      & 58 & 0.7545 & 92.89 &  98.3 & 3.35 & 0 \\
crambin (46 aa)   & 46 & 0.5816 & 83.31 &  89.1 & 3.33 & 0 \\
\bottomrule
\end{tabular}
\end{center}

All three pass the gate, so the harness folds, and the choice of $3.0\ang$
as the clash line is borne out empirically (confident folds have their
closest pair at $3.3$--$3.7\ang$).

\subsection{HA constructs}

\begin{center}
\small
\setlength{\tabcolsep}{4pt}
\begin{tabular}{@{}lrrrrr@{}}
\toprule
 & cleaved & mature & HA1 alone & HA1 alone & HA2 alone \\
\midrule
positions / steps & 551 / 256 & 550 / 256 & 329 / 256 & 329 / 329 & 221 / 221$^{*}$ \\
pTM & 0.2322 & 0.2614 & 0.4642 & 0.4133 & 0.5674 \\
pLDDT, all residues & 63.89 & 59.20 & 65.16 & 59.83 & 93.93 \\
\% residues $\ge 70$ & 45.5 & 31.8 & 37.1 & 18.5 & 97.7 \\
pLDDT, HA1 (\% $\ge 70$) & 63.22 (38.6) & 51.42 (11.2) & 65.16 (37.1) & 59.83 (18.5) & --- \\
pLDDT, HA2 (\% $\ge 70$) & 64.90 (55.7) & 70.78 (62.4) & --- & --- & 93.93 (97.7) \\
min C$\beta$--C$\beta$ ($\ang$) & 1.50 & 1.17 & 2.01 & 0.44 & 4.14 \\
clashes & 15 & 9 & 7 & 16 & 0 \\
gate & fail & fail & fail & fail & \textbf{pass} \\
Cys14--Cys137 C$\beta$ ($\ang$) & 95.67 & 129.74 & --- & --- & --- \\
pLDDT at Cys14 / Cys137 & 73.6 / 61.3 & 22.5 / 80.9 & --- & --- & --- \\
\bottomrule
\end{tabular}
\end{center}

\noindent The first run, at one decoding step, gave pTM 0.1827 and a bridge
distance of $38.55\ang$ with a $1.64\ang$ closest cysteine pair. It is not
tabulated because it measured a structure the model did not believe in;
see Section~\ref{sec:corrections}. $^{*}$HA2 alone was requested at 256 steps
and clamped to 221 by the script (a chain cannot have more steps than
positions), so it has been seen at one step count only.

\paragraph{Reproducibility.} Cleaved and mature were each run twice, with
different script versions, on the same machine, and gave identical pTM
(0.2322 and 0.2614) and identical bridge distances (95.67 and $129.74\ang$).
HA2 alone was run twice as well (the second time by the clamp above) with
identical output to every printed digit. Decoding at temperature zero looks
deterministic here. This is three repeat pairs on one machine, not a proof.

\subsection{Step sweep}

\begin{center}
\small
\setlength{\tabcolsep}{6pt}
\begin{tabular}{@{}lrrr@{}}
\toprule
 & 64 steps & 128 steps & 256 steps \\
\midrule
pLDDT, HA1 in cleaved   & 56.16 & 64.32 & 63.22 \\
pLDDT, HA1 fused        & 61.35 & 61.90 & 51.42 \\
pLDDT, HA1 alone        & 48.43 & 67.01 & 65.16 \\
pLDDT, HA2 in cleaved   & 68.49 & 63.27 & 64.90 \\
pLDDT, HA2 fused        & 68.94 & 71.40 & 70.78 \\
pLDDT, HA2 alone$^{*}$  & 93.67 & 88.47 & 93.93 \\
pLDDT at Cys14, cleaved & 31.7 & 58.0 & 73.6 \\
pLDDT at Cys14, fused   & 41.1 & 69.9 & 22.5 \\
clashes, cleaved / fused & 22 / 9 & 19 / 2 & 15 / 9 \\
clashes, HA1 alone / HA2 alone & 17 / 0 & 5 / 0 & 7 / 0 \\
bridge C$\beta$--C$\beta$, cleaved ($\ang$) & 98.8 & 137.9 & 95.7 \\
bridge C$\beta$--C$\beta$, fused ($\ang$)   & 164.3 & 48.0 & 129.7 \\
\bottomrule
\end{tabular}

\smallskip
\begin{tabular}{@{}lrrrl@{}}
\toprule
Contrast at equal steps & 64 & 128 & 256 & sign \\
\midrule
HA1 pLDDT: break $-$ fused  & $-5.2$ & $+2.4$ & $+11.8$ & mixed \\
HA1 pLDDT: alone $-$ cleaved & $-7.7$ & $+2.7$ & $+1.9$ & mixed \\
Cys14 pLDDT: break $-$ fused & $-9.4$ & $-11.9$ & $+51.1$ & mixed \\
HA2 pLDDT: alone $-$ cleaved & $+25.2$ & $+25.2$ & $+29.0$ & $+$ at all \\
HA2 pLDDT: alone $-$ fused   & $+24.7$ & $+17.1$ & $+23.2$ & $+$ at all \\
HA2 pLDDT: break $-$ fused   & $-0.5$ & $-8.1$ & $-5.9$ & $-$ at all \\
\bottomrule
\end{tabular}
\end{center}
\noindent $^{*}$HA2 alone is capped at its 221 residues, so its 256-step
column is the 221-step run. Computed by \path{bridge_test/sweep_contrasts.py};
$n=3$, so means and spreads are descriptive and no p-value is claimed.

\subsection{Findings}

\paragraph{The chain-break's benefit to HA1 is a 256-step effect.} Break minus
fused is $-5.2$, $+2.4$, $+11.8$ at 64, 128, 256 steps. The $+11.8$ comes from
fused HA1 dipping to 51.42 at 256 steps (61.35 and 61.90 at the other two),
while cleaved HA1 stays at 56--64. The Cys14 anchor result is the same
artefact: at 64 and 128 steps the break gives the \emph{lower} anchor
confidence. Neither claim passes the sign-at-every-step rule.

\paragraph{The clash \emph{topology} differs.} At 256 steps mature has fewer
clashes (9 against 15) but they sit at $|i-j|$ of 80--159, so distant segments
are superimposed; cleaved's are local ($|i-j|\le 16$). Across the sweep this
holds at two of three step counts (fused is long-range at 128 and 256 but not
at 64; cleaved has one long-range clash, at 182, at 128 steps), so it is
suggestive, not clean.

\paragraph{HA2 is far more confident alone than in any assembly.} HA2 alone
minus HA2 in the cleaved complex is $+25.2$, $+25.2$, $+29.0$, and against the
fused one $+24.7$, $+17.1$, $+23.2$: the same sign at every step count, with
spreads of 2--4 points on differences of 17--29. This is the one robust
contrast in the study. The break costs HA2 a little ($-0.5$, $-8.1$, $-5.9$,
negative throughout, but of the same size as the step noise).

\paragraph{HA1 alone against HA1 in the complex is undetermined:} $-7.7$,
$+2.7$, $+1.9$. At 64 steps HA1 alone is unfolded (0.0\% of residues above
70). Whether the complex is right in a way pLDDT cannot see, for example in
the interchain placement, is untested; the assembled pTM is too low to
resolve it.

\paragraph{Step count is a large, non-monotone effect.} HA1 alone scores
48.43, 67.01, 65.16 and 59.83 at 64, 128, 256 and 329 steps; it rises,
plateaus and falls. Fused HA1 drops 10 points between 128 and 256 steps while
cleaved HA1 does not. Which cysteines pair also moves with steps alone
(one:64--76 goes from $3.92\ang$ to $26.64\ang$ between 329 and 256 steps).

\paragraph{HA1 is the hard half in every construct.} Alone it fails the gate
at both step counts tried: 16 clashes at 329 steps (minimum C$\beta$
separation $0.44\ang$) and 7 at 256 steps ($2.01\ang$); it fails at 64 and 128 too.

\paragraph{The bridge is not formed, at any step count.} Cleaved gives
$98.8$, $137.9$, $95.7\ang$ and fused $164.3$, $48.0$, $129.7\ang$, against a
$4.5\ang$ criterion. No assembled run passes the gate either (cleaved 22, 19,
15 clashes; fused 9, 2, 9, the 2-clash run failing on pTM), and HA1 alone fails
at every step count tried.

\subsection{Against experimental structures}

Three experimental entries were fetched, and each was accepted only if its own
sequence matched the reference: BPTI (5PTI, $1.00\ang$), crambin (3NIR,
$0.48\ang$) and protein~G (1PGB, $1.92\ang$). C$\beta$ is taken by the model's
own convention on both sides; the deposited C$\beta$ differs by about
$0.1$--$0.2\ang$.

\begin{center}
\small
\setlength{\tabcolsep}{5pt}
\begin{tabular}{@{}lrrr@{}}
\toprule
 & protein G & BPTI & crambin \\
\midrule
CA RMSD ($\ang$) & 2.96 & 1.22 & 7.91 \\
TM-score (approx.) & 0.835 & 0.922 & 0.316 \\
contact F1, $|i-j|\ge 6$ & 0.80 & 0.92 & 0.35 \\
local CA--CA error, all pairs ($\ang$) & 1.10 & 0.41 & 4.06 \\
worst 7-residue window ($\ang$) & 2.68 (11--17) & 0.78 (52--58) & 3.41 (12--18) \\
pTM / mean pLDDT (model) & 0.76 / 93 & 0.75 / 93 & 0.58 / 83 \\
\bottomrule
\end{tabular}
\end{center}

\paragraph{The crambin fold is wrong, and passes the gate.} A TM-score of $0.32$
is well below the $0.5$ that marks a shared fold, and the local error is $4\ang$
on average. The model gives it pTM $0.58$ and pLDDT $83$, with no clashes. The
geometry gate is necessary, not sufficient.

\paragraph{Protein~G's fold is right, and one loop is confidently wrong.}
Residues 9--15 lie $5$--$11\ang$ from experiment at pLDDT $80$--$91$. Their
internal shape is wrong too: the intra-loop CA--CA error averages $1.8\ang$, and
the end-to-end distance is $15.3\ang$ against $11.3\ang$. Their experimental
B-factors are normal, so this is not crystal disorder, and confidence does not
flag the worst residues.

\paragraph{The $4.5\ang$ criterion is exact on experimental geometry.} All six
supplied pairings are disulfides in the files (SG--SG $2.02$--$2.05\ang$). Exactly
the native bonds lie within $4.5\ang$ in both proteins. The nearest non-bonded
pair is $6.14\ang$ in BPTI and $5.79\ang$ in crambin. The earlier claim that the
criterion lacks specificity was wrong: the non-native crambin pair Cys16--Cys40
is predicted at $3.33\ang$ and is $17.9\ang$ in the experiment. That specificity
failure belongs to the prediction.

\begin{center}
\small
\begin{tabular}{@{}llrrrl@{}}
\toprule
Protein & Pair & SG--SG ($\ang$) & exp.\ C$\beta$ & pred.\ C$\beta$ & pred.\ within $4.5\ang$ \\
\midrule
BPTI    & Cys5--Cys55  & 2.04 & 3.49 & 3.66 & yes \\
BPTI    & Cys14--Cys38 & 2.03 & 3.89 & 3.61 & yes \\
BPTI    & Cys30--Cys51 & 2.02 & 3.71 & 4.45 & yes \\
crambin & Cys3--Cys40  & 2.03 & 3.45 & 7.75 & no \\
crambin & Cys4--Cys32  & 2.05 & 4.03 & 3.79 & yes \\
crambin & Cys16--Cys26 & 2.04 & 3.82 & 10.42 & no \\
\midrule
crambin & Cys16--Cys40 (non-native) & -- & 17.88 & 3.33 & \textbf{yes} \\
\bottomrule
\end{tabular}
\end{center}

\paragraph{The round trip holds both folds.} Experimental coordinates encoded to
structure tokens and decoded reproduce the native fold to $0.31\ang$ (BPTI) and
$0.30\ang$ (crambin), with TM-scores $0.99$ and $0.98$. The known bonds come back
within about $0.4\ang$. The representation can hold these folds and these bonds.
For crambin the failure is in generation, not in the representation.

\paragraph{BPTI cannot test imposition.} Its bonds form without any prompt. A
prompt on Cys14 and Cys38 leaves all three bond distances within $0.02\ang$ of
the unprompted values.

\paragraph{A coordinate prompt pulls crambin's pair partway, and costs confidence.}
Prompting the backbone of residues 3 and 40 moves their CA--CA distance from
$7.84$ to $6.68\ang$ (experiment $5.66\ang$; $53\%$ of the gap closed) and their
C$\beta$--C$\beta$ distance from $7.75$ to $5.14\ang$ (experiment $3.45\ang$;
$61\%$ closed). It forms no other bond, the fold stays wrong (RMSD $7.9$ to
$7.1\ang$), and pTM falls from $0.58$ to $0.46$. This is a soft, partial, local
pull, not an imposition. The crambin fold is wrong, so the test is confounded
either way.

% ===========================================================================
\section{Corrections made during the session}
\label{sec:corrections}
% ===========================================================================

These are recorded because several first readings were wrong.

\begin{enumerate}
\item \textbf{One decoding step.} The first verdict (``bridge not formed,
  $38.55\ang$'') came from a script that used the default
  \texttt{num\_steps=1}, which samples all 551 structure tokens at once. It
  scored pTM 0.18 with a $1.64\ang$ ``contact'' shorter than a C--C bond.
  The verdict was worthless. The gate now runs before any verdict.
\item \textbf{A control that did not control.} Protein~G validated the fold
  but has no cysteines, so it could not validate the measurement. BPTI and
  crambin were added for that reason.
\item \textbf{Two overstated claims, retracted.} The calibration script
  printed ``$4.5\ang$ is sound; HA2's disulfides are not formed'' from BPTI
  alone. That tested sensitivity on one protein and specificity not at all.
  Separately, HA1-alone pairs at $3.51$ and $3.92\ang$ had been described as
  ``real disulfide geometry''; by the crambin result they could be
  non-native.
\item \textbf{A pre-registered prediction failed.} I expected the trimer
  signature (both chains confident, assembly low). The result was
  asymmetric instead: HA2 adequate, HA1 poor.
\item \textbf{``HA2's presence helps HA1'' retracted.} It rested on HA1
  alone at 329 steps against HA1 in the complex at 256. At matched steps HA1
  alone is 65.16 against 63.22. This report had listed that confound as its
  first limitation, and the rerun resolved it against the claim. I had also
  said the premise ``ESM3 should not treat the chains as separate'' was
  confirmed; what is confirmed is narrower (an unbroken fusion is wrong, and
  the break repairs it).
\item \textbf{Three more findings withdrawn after the step sweep.} The
  chain-break benefit to HA1, its benefit to the Cys14 anchor, and the
  narrower ``break repairs fusion'' that replaced the first claim all rested
  on 256 steps. I stated each with more confidence than one step count could
  support. The sweep also exposed a bug in my own contrast script, which kept
  an old pre-pLDDT result file for the 256-step runs; an empty column gave it
  away.
\item \textbf{Two calibration verdicts, retracted.} The crambin calibration
  called the $4.5\ang$ criterion unreliable, but it was computed on a misfolded
  protein and says nothing about the criterion. The claim that the criterion
  lacks specificity was also wrong: experimental geometry shows it exact, and
  the non-native crambin pair is a model error. The supplied disulfide pairings
  are now confirmed as bonds in the experimental files.
\item \textbf{Small bugs.} The chain-break position was admitted into the
  C$\beta$ statistics; labels were wrong for the break-free reference;
  a results file was overwritten by a generic filename.
\end{enumerate}

% ===========================================================================
\section{Limitations}
% ===========================================================================

\subsection{Of the measurements}

\begin{itemize}
\item \textbf{Experimental ground truth exists only for the controls.} There is
  no experimental HA structure, so nothing about HA's fold or bridge has been
  compared with experiment. Each control was predicted at one step count.
\item \textbf{C$\beta$--C$\beta$ is a proxy.} The decoder emits no sulfur, and
  C$\beta$ is inferred from the backbone. A bond criterion built on it
  inherits any backbone error, and glycines (which have no C$\beta$) are
  invisible to the clash count.
\item \textbf{Three step counts per construct.} Enough to show that the sign
  of several contrasts flips, not enough to characterise the step
  dependence. HA2 alone is capped at 221 steps. It is unknown whether steps or
  tokens decoded per step is the operative variable.
\item \textbf{One run per condition.} The three repeat pairs agree exactly, but
  no condition has been run across seeds or checkpoints, and nothing here
  gives error bars.
\item \textbf{pTM depends on length,} so comparing it between a 46-residue
  and a 550-residue chain is weak. Only same-construct comparisons are
  clean.
\item \textbf{Two controls.} BPTI is small and folds correctly, so it cannot test
  imposition. Crambin's fold is wrong, so its results are confounded. Neither
  says what the criterion does in a 550-residue fold.
\item \textbf{The control pairings are confirmed by experiment.} Each supplied
  pairing is a disulfide in the deposited file (SG--SG $2.02$--$2.05\ang$). The
  sequences were checked against the experimental SEQRES, not against UniProt.
\item \textbf{Only \texttt{esm3-open}.} No other checkpoint was tried.
\end{itemize}

\subsection{Of the scope}

\begin{itemize}
\item \textbf{Monomer only.} HA is a trimer. A single HA1+HA2 protomer may
  have no self-contained fold, and a low score may then be correct. Not
  tested.
\item \textbf{No verified native disulfide list for HA,} so partner recovery
  cannot be run on HA at all. Positive bridge claims about HA are not
  possible yet.
\item \textbf{The energy machinery cannot take a chain break.} The sampler's
  soft-probability path has no way to carry token 31, so none of the
  \(U_{\mathrm{am}}\) work has been applied to the cleaved form.
\item \textbf{Glycans, the fusion peptide, and everything outside the
  sequence} are absent by construction.
\end{itemize}

\subsection{Of the session itself}

\begin{itemize}
\item For most of the session I could not run code: the shell was blocked
  until auto mode was switched off in this final turn. Every script was
  therefore written unexecuted and validated by your runs. One consequence:
  \texttt{test\_cleaved\_complex.py} still carries an unused
  \texttt{present\_atoms()} helper that I reintroduced by mistake.
\item One commit message (\texttt{87dc571}) has a stray line that refers to
  non-existent session notes. The file content is fine; the message cannot be
  amended without rewriting history.
\end{itemize}

% ===========================================================================
\section{To do}
% ===========================================================================

\subsection{Recommended next step}

Run the control step sweep. Each control was predicted at one step count, and
the HA sweep showed that step count can change a fold by more than the effects
under test. Crambin's wrong fold is a property of the protein only if it is
wrong at 11, 23 and 46 steps. If it is right at some step count, the crambin
calibration has to be redone at that count. The sweep is
\path{run_pdb_comparison.sh controls}, and its reading was fixed in advance: a
fold is correct if the CA RMSD is at most $2.5\ang$ and the TM-score at least
$0.7$, and a control is reliable only if correct at every step count of at
least $L/4$.

\subsection{Prioritised list}

\begin{enumerate}
\item \textbf{Control step sweep} (above).
\item \textbf{A local conformation check} for protein~G's loop. pLDDT does not
  flag residues 9--15, so confidence is not an error check; a local geometry
  test against experiment is what separates the confident errors.
\item \textbf{One verdict table for experiment and prediction,} as in the
  experimental comparison, stating the fold's correctness before any distance.
\item \textbf{Interface contacts, cleaved against fused,} once the contact
  metric is validated against experiment on the short proteins.
\item \textbf{A verified native disulfide list for HA,} from a named
  structure entry, asserted at import the way the controls are. Required
  before any partner-recovery claim about HA.
\item \textbf{The trimer:} \texttt{one|two|one|two|one|two}, 1653 positions.
  The biologically correct question, and expensive.
\item \textbf{Condition on a real HA structure} and ask whether ESM3
  \emph{maintains} the bridge rather than discovers it. A different and more
  tractable question.
\item \textbf{Only if the cleaved form is needed in the energy machinery:}
  inject a one-hot at token 31 in the padded region, set the structure token
  at the break to \texttt{STRUCTURE\_CHAINBREAK\_TOKEN} in
  \texttt{\_default()}, and exclude the break site from \texttt{mutate()}.
\end{enumerate}

\subsection{Housekeeping}

\begin{itemize}
\item At session end, append the two devlogs to the main one:
\begin{verbatim}
cat bridge_test/DEVLOG_bridge_test.txt \
    bridge_test/DEVLOG_calibration.txt >> DEVLOG.txt
\end{verbatim}
\item Give results files run-identifying names; the generic default
  overwrote the protein~G control once.
\item Remove the dead \texttt{present\_atoms()} helper.
\item \path{bridge_test/summarize_runs.py} tabulates every result file from
  its content; use it instead of copying numbers by hand.
\end{itemize}

% ===========================================================================
\section{Reproducing the runs}
% ===========================================================================

From \path{softmax/code}:

\begin{verbatim}
python tests/test_cleaved_complex.py --ref protein_g --steps 56
python tests/test_cleaved_complex.py --ref bpti      --steps 58
python tests/test_cleaved_complex.py --ref crambin   --steps 46
python tests/test_cleaved_complex.py --ref cleaved   --steps 256
python tests/test_cleaved_complex.py --ref mature    --steps 256
python tests/test_cleaved_complex.py --ref one       --steps 329
python tests/test_cleaved_complex.py --ref one       --steps 256
python tests/test_cleaved_complex.py --ref two       --steps 221
\end{verbatim}

The experimental comparison is \path{bridge_test/run_pdb_comparison.sh}, with
phases \texttt{fetch}, \texttt{predict}, \texttt{compare}, \texttt{controls},
\texttt{roundtrip} and \texttt{prompt}. Output files are in
\path{softmax/code/bridge_test/}.

\end{document}
